\documentclass[11pt,a4paper]{article}

% --- Codificación e idioma ---
\usepackage[utf8]{inputenc}
\usepackage[T1]{fontenc}
\usepackage[spanish,es-tabla]{babel}
\usepackage{lmodern}

% --- Geometría y formato general ---
\usepackage[margin=2.5cm]{geometry}
\usepackage{microtype}
\usepackage{parskip}
\usepackage{enumitem}
\setlist{topsep=2pt,itemsep=2pt,parsep=0pt}

% --- Tablas ---
\usepackage{array}
\usepackage{booktabs}
\usepackage{longtable}

% --- Cabeceras y numeración ---
\usepackage{fancyhdr}
\pagestyle{fancy}
\fancyhf{}
\fancyhead[L]{\small Bases de Datos II}
\fancyhead[R]{\small Resumen Teórico --- Práctico N.\textsuperscript{o} 1}
\fancyfoot[C]{\thepage}
\renewcommand{\headrulewidth}{0.4pt}

% --- Color y código ---
\usepackage{xcolor}
\definecolor{codebg}{RGB}{248,248,248}
\definecolor{codeframe}{RGB}{200,200,200}
\definecolor{codekw}{RGB}{0,0,180}
\definecolor{codestr}{RGB}{163,21,21}
\definecolor{codecmt}{RGB}{0,128,0}
\definecolor{notebg}{RGB}{255,250,230}
\definecolor{noteframe}{RGB}{220,180,80}

\usepackage{listings}
\lstdefinelanguage{SQLext}[]{SQL}{
    morekeywords={NUMERIC,VARCHAR,DECIMAL,DATE,TIMESTAMP,TIME,
        BIGINT,VARCHAR2,NUMBER,ENUM,SCHEMA,SEQUENCE,TRIGGER,
        BEFORE,EACH,ROW,NEXTVAL,DUAL,RAISE,EXCEPTION,LOOP,
        IDENTIFIED,QUOTA,UNLIMITED,USERS,TABLESPACE,
        DATABASE,DOMAIN,GRANT,REVOKE,ROLE,USER,LOGIN,NOLOGIN,
        OPTION,REFERENCES,RESTRICT,CASCADE,
        AUTO\_INCREMENT,IF,EXISTS,ALTER,SHOW,USE,
        SEARCH\_PATH,EXECUTE,IMMEDIATE},
    sensitive=false,
    morecomment=[l]{--},
    morestring=[b]'
}
\lstset{
    language=SQLext,
    basicstyle=\ttfamily\footnotesize,
    keywordstyle=\color{codekw}\bfseries,
    stringstyle=\color{codestr},
    commentstyle=\color{codecmt}\itshape,
    backgroundcolor=\color{codebg},
    frame=single,
    rulecolor=\color{codeframe},
    breaklines=true,
    breakatwhitespace=true,
    columns=fullflexible,
    keepspaces=true,
    showstringspaces=false,
    captionpos=b,
    xleftmargin=0.4em,
    xrightmargin=0.4em,
    aboveskip=8pt,
    belowskip=8pt,
    tabsize=2
}

% --- Cajas de nota ---
\usepackage[most]{tcolorbox}
\newtcolorbox{nota}{colback=notebg,colframe=noteframe,arc=2pt,
    left=6pt,right=6pt,top=4pt,bottom=4pt,boxrule=0.6pt}

% --- Hipervínculos ---
\usepackage[hidelinks,bookmarks=true]{hyperref}

\title{Resumen Teórico --- Práctico N.\textsuperscript{o} 1: Repaso de SQL\\[4pt]
       \large Bases de Datos II}
\author{Tomás Rodeghiero}
\date{2026}

\begin{document}

% =====================================================
% Carátula
% =====================================================
\begin{titlepage}
    \centering
    \vspace*{1.5cm}

    {\large \textbf{Universidad Nacional de Río Cuarto}}\\[6pt]
    {\normalsize Facultad de Ciencias Exactas, Físico-Químicas y Naturales}\\[4pt]
    {\normalsize Departamento de Computación}\\[4pt]
    {\normalsize Licenciatura en Ciencias de la Computación}\\[3cm]

    {\Large \textbf{Bases de Datos II}}\\[1.2cm]

    {\Large \textbf{Resumen Teórico}}\\[6pt]
    {\large Práctico N.\textsuperscript{o} 1 --- Repaso de SQL}\\[3cm]

    \begin{tabular}{rl}
        \textbf{Alumno:} & Tomás Rodeghiero \\[4pt]
        \textbf{Año:}    & 2026 \\
    \end{tabular}

    \vfill
    {\small Material elaborado a partir de los teóricos
    \emph{Teorico\_1\_Introduccion} y \emph{Teorico\_2\_Repaso\_SQL}
    de la cátedra y de las resoluciones del práctico.}
\end{titlepage}

\tableofcontents
\newpage

% =====================================================
\section{Introducción}
% =====================================================

Bases de Datos II retoma los conceptos relacionales de Bases de Datos I y los
profundiza hacia la administración real de motores: optimización, transacciones,
concurrencia, programación en el servidor (procedimientos almacenados,
\emph{triggers}), \mbox{XML/JSON}, bases NoSQL, bases geográficas e inteligencia
de negocios. El primer práctico es un \textbf{repaso integral de SQL} que se
apoya en los dos primeros teóricos de la materia y deja preparada la base
sobre la cual se construirán los temas posteriores.

El práctico tiene tres focos claros: definir esquemas con DDL respetando las
restricciones del enunciado, poblarlos con datos coherentes mediante DML y
escribir consultas combinando filtros, \emph{joins}, agrupamiento y agregación.
Como detalle no menor, el práctico se resuelve sobre los tres motores que la
cátedra exige manejar: \textbf{MySQL 8+}, \textbf{PostgreSQL 14+} y
\textbf{Oracle 10/11 XE}.

\begin{nota}
\textbf{Idea clave.} Aunque el SQL estándar es uno solo, cada motor tiene
diferencias prácticas importantes (autoincrement, tipos numéricos, dominios,
manejo de esquemas, etc.). Aprender a leer SQL en el estándar y traducirlo
a cada motor es uno de los objetivos centrales del práctico.
\end{nota}

% =====================================================
\section{SQL: visión general}
% =====================================================

\textbf{SQL} (\emph{Structured Query Language}) es el lenguaje estándar para
organizar, gestionar y recuperar datos en bases de datos relacionales.
Históricamente nació como \emph{SEQUEL} en un prototipo de IBM y fue luego
estandarizado por ANSI e ISO. Con el tiempo el estándar incorporó
\emph{triggers} (SQL:1999), interacción con XML (SQL:2005), funciones
polimórficas, JSON y funciones trigonométricas (SQL:2016), entre otras
extensiones. En paralelo, los motores comerciales (Oracle, SQL Server,
PostgreSQL, MySQL) ofrecen sus propias variantes y agregan funcionalidades
fuera del estándar.

SQL se divide tradicionalmente en tres sublenguajes:

\begin{description}[leftmargin=2.2cm,style=nextline]
    \item[\textbf{DML} (Data Manipulation Language)]
    Manipulación de datos: consultas (\texttt{SELECT}) basadas en álgebra
    relacional, e inserción, modificación y borrado de filas
    (\texttt{INSERT}, \texttt{UPDATE}, \texttt{DELETE}).

    \item[\textbf{DDL} (Data Definition Language)]
    Definición de la estructura: creación, modificación y borrado de tablas,
    dominios, índices, vistas y demás objetos del esquema
    (\texttt{CREATE}, \texttt{ALTER}, \texttt{DROP}).

    \item[\textbf{DCL} (Data Control Language)]
    Control de acceso y permisos sobre los objetos
    (\texttt{GRANT}, \texttt{REVOKE}). Es el foco del Práctico 2,
    pero conviene reconocerlo desde ahora para ubicar cada sentencia
    en su categoría.
\end{description}

El Práctico 1 trabaja casi exclusivamente con \textbf{DDL} (Ejercicio 1)
y \textbf{DML} (Ejercicios 2 y 3).

% =====================================================
\section{DML: manipulación de datos}
% =====================================================

% -----------------------------------------------------
\subsection{La sentencia \texttt{SELECT}}
% -----------------------------------------------------

\texttt{SELECT} es el comando más usado de SQL. Recupera datos de una o varias
tablas y se organiza en seis cláusulas, dos obligatorias y cuatro opcionales:

\begin{lstlisting}
SELECT [ALL|DISTINCT] { * | expr_columna_1 [AS alias_1] [, ...] }
  FROM nombre_tabla_1 [alias_1] [, nombre_tabla_n [alias_n]]
 [WHERE  condicionWhere]
 [GROUP BY expr_columna_group1 [, expr_columna_group2] [...]]
 [HAVING  condicionHaving]
 [{UNION [ALL] | INTERSECT | EXCEPT} SELECT ...]
 [ORDER BY nombre_campo_1 [ASC|DESC] [, nombre_campo_2 [ASC|DESC]]];
\end{lstlisting}

Conceptualmente, el orden \emph{lógico} en que el motor evalúa una consulta es
distinto al orden en que se escribe: primero \texttt{FROM} (con sus
\texttt{JOIN}s), luego \texttt{WHERE}, después \texttt{GROUP BY}, luego
\texttt{HAVING}, después la lista del \texttt{SELECT} y finalmente
\texttt{ORDER BY}. Tener este modelo mental en la cabeza explica, por ejemplo,
por qué un alias definido en el \texttt{SELECT} no siempre puede usarse
en el \texttt{WHERE}.

\textbf{\texttt{DISTINCT} y \texttt{ALL}.} \texttt{DISTINCT} elimina filas
repetidas del resultado. \texttt{ALL} (valor por defecto) las conserva. La
diferencia es relevante cuando proyectamos sólo unas pocas columnas y muchas
combinaciones se repiten.

\textbf{Alias.} Tanto las tablas como las columnas pueden recibir alias con
\texttt{AS}. Es habitual abreviar las tablas (\texttt{cliente c},
\texttt{factura f}) para escribir condiciones de \emph{join} más cortas.

% -----------------------------------------------------
\subsection{Cláusula \texttt{WHERE}: filtros sobre filas}
% -----------------------------------------------------

\texttt{WHERE} se usa para quedarse sólo con las filas que cumplan una
condición lógica. Una condición es una \emph{expresión}: compara columnas
contra valores o contra otras columnas y combina sub-expresiones con
operadores lógicos.

\paragraph{Operadores lógicos.}
\texttt{AND}, \texttt{OR}, \texttt{NOT} y \texttt{XOR}. Como en cualquier
lenguaje, \texttt{AND} tiene mayor precedencia que \texttt{OR}, así que
conviene siempre parentizar para evitar ambigüedades.

\paragraph{Operadores de comparación.}

\begin{itemize}
    \item Comparadores: \texttt{=}, \texttt{<>}, \texttt{<}, \texttt{>},
    \texttt{<=}, \texttt{>=}.
    \item \texttt{BETWEEN a AND b}: pertenece al intervalo cerrado $[a,b]$;
    aplica a tipos numéricos, fechas y cadenas.
    \item \texttt{LIKE 'patrón'}: comparación por patrón. \texttt{\%}
    representa cualquier secuencia de caracteres y \texttt{\_} un único
    carácter.
    \item \texttt{IN (subconsulta o lista)}: verdadero si el valor está dentro
    del conjunto. Muy usado para conectar consultas anidadas.
    \item \texttt{EXISTS (subconsulta)}: verdadero si la subconsulta devuelve
    al menos una fila. Particularmente útil para subconsultas correlacionadas.
    \item \texttt{IS NULL} / \texttt{IS NOT NULL}: el único modo correcto de
    comparar contra \texttt{NULL}, ya que \texttt{= NULL} \emph{nunca} es
    verdadero (la comparación da \texttt{UNKNOWN}).
\end{itemize}

\paragraph{Ejemplos sobre el esquema del teórico.}

\begin{lstlisting}
-- Artículos cuya descripción empieza con A.
SELECT * FROM articulos WHERE descr LIKE 'A%';

-- Empieza con A y termina con K.
SELECT * FROM articulos WHERE descr LIKE 'A%' AND descr LIKE '%K';

-- Artículos provistos por algun proveedor (subconsulta con IN).
SELECT * FROM articulos
 WHERE Nart IN (SELECT Nart FROM provee);
\end{lstlisting}

% -----------------------------------------------------
\subsection{Cláusula \texttt{ORDER BY}: ordenamiento}
% -----------------------------------------------------

\texttt{ORDER BY} ordena el resultado de la consulta por una o más columnas.
Sin esta cláusula, el orden devuelto por el motor no está garantizado, aunque
suele coincidir con el primer índice utilizado.

\begin{lstlisting}
ORDER BY expresion_orden1 [ASC|DESC]
       [, expresion_orden2 [ASC|DESC]] [, ...];
\end{lstlisting}

La expresión puede ser un nombre de columna, una expresión cualquiera o el
\emph{número de posición} de la columna en el \texttt{SELECT}. \texttt{ASC}
es ascendente (por defecto) y \texttt{DESC} descendente. Si se ordena por
varias columnas, la segunda actúa como criterio de desempate dentro de la
primera, y así sucesivamente.

\paragraph{Ejemplo.}

\begin{lstlisting}
-- Listado de articulo - proveedor - precio, ordenado por Nart ascendente.
SELECT DISTINCT Nart AS numero_art,
                pv.Nprov AS numero_prov,
                nombre,
                precio_venta
  FROM Proveedores pd, Provee pv
 WHERE pd.Nprov = pv.Nprov
 ORDER BY Nart;
\end{lstlisting}

% -----------------------------------------------------
\subsection{Funciones de agregación}
% -----------------------------------------------------

Las \emph{funciones agregadas} reciben una expresión de columna y devuelven
un único valor calculado sobre todas las filas. Las del estándar son:

\begin{center}
\begin{tabular}{@{}l p{10cm}@{}}
\toprule
\textbf{Función} & \textbf{Devuelve} \\
\midrule
\texttt{AVG(expr)}   & Promedio de los valores no nulos. \\
\texttt{COUNT(expr)} & Cantidad de valores no nulos. \texttt{COUNT(*)}
                       cuenta filas, incluyendo aquéllas con \texttt{NULL}. \\
\texttt{SUM(expr)}   & Suma de los valores no nulos. \\
\texttt{MIN(expr)}   & Mínimo (compatible con números, fechas y cadenas). \\
\texttt{MAX(expr)}   & Máximo. \\
\bottomrule
\end{tabular}
\end{center}

\begin{nota}
Salvo \texttt{COUNT(*)}, las funciones agregadas \textbf{ignoran los valores
\texttt{NULL}}. Esto es importante a la hora de calcular promedios: el
denominador es la cantidad de valores presentes, no la cantidad de filas.
\end{nota}

\paragraph{Ejemplo simple.}

\begin{lstlisting}
SELECT AVG(precio) AS promedio FROM articulos;
\end{lstlisting}

% -----------------------------------------------------
\subsection{\texttt{GROUP BY}: consultas sumarias}
% -----------------------------------------------------

\texttt{GROUP BY} convierte un \texttt{SELECT} en una \emph{consulta sumaria}:
en lugar de devolver una fila por cada fila de la base, agrupa todas las que
comparten el valor de las columnas indicadas y produce una fila por grupo.
Las columnas que aparecen en el \texttt{SELECT} y no están dentro de una
función de agregación deben aparecer también en el \texttt{GROUP BY}.

\paragraph{Ejemplo.}

\begin{lstlisting}
-- Para cada articulo, precio promedio entre todos sus proveedores.
SELECT Nart, AVG(precio_venta) AS promedio_p_v
  FROM Provee
 GROUP BY Nart;
\end{lstlisting}

% -----------------------------------------------------
\subsection{\texttt{HAVING}: filtros sobre grupos}
% -----------------------------------------------------

\texttt{HAVING} es a los grupos lo que \texttt{WHERE} es a las filas: filtra
los grupos producidos por \texttt{GROUP BY} usando condiciones que pueden
involucrar funciones agregadas. La regla práctica es:

\begin{itemize}
    \item Si la condición sólo usa \emph{atributos planos} (no agregados),
    debe ir en \texttt{WHERE} (más eficiente: filtra antes de agrupar).
    \item Si la condición involucra \emph{funciones agregadas}, debe ir
    en \texttt{HAVING} (sólo se puede evaluar después de armar los grupos).
\end{itemize}

\paragraph{Ejemplo.}

\begin{lstlisting}
-- Articulos cuyo precio promedio de venta entre proveedores es < 1000.
SELECT Nart, AVG(precio_venta)
  FROM Provee
 GROUP BY Nart
HAVING AVG(precio_venta) < 1000;
\end{lstlisting}

% -----------------------------------------------------
\subsection{Subconsultas}
% -----------------------------------------------------

Una \emph{subconsulta} es un \texttt{SELECT} usado dentro de otra sentencia.
Pueden aparecer en distintos lugares y con distintas semánticas:

\begin{itemize}
    \item En el \texttt{WHERE} junto con \texttt{IN}, \texttt{NOT IN},
    \texttt{EXISTS}, \texttt{NOT EXISTS} u operadores de comparación
    (\texttt{=}, \texttt{<}, \texttt{ANY}, \texttt{ALL}).
    \item En el \texttt{FROM} como \emph{tabla derivada}, dándole un alias
    para tratarla como una relación más.
    \item En el \texttt{SELECT} como subconsulta escalar, devolviendo un
    único valor por fila.
\end{itemize}

\textbf{Subconsultas correlacionadas.} Son aquellas en las que la subconsulta
hace referencia a una columna de la consulta exterior. Se ejecutan
conceptualmente una vez por cada fila externa. Aparecen típicamente en
patrones del tipo ``existe alguien con \dots'' o ``no existe ninguno
con \dots''.

\paragraph{Ejemplo de \texttt{NOT EXISTS}.}

\begin{lstlisting}
-- Clientes que no tienen ninguna factura asociada.
SELECT *
  FROM cliente c
 WHERE NOT EXISTS (SELECT 1
                     FROM factura f
                    WHERE f.nro_cliente = c.nro_cliente)
 ORDER BY c.apellido DESC, c.nombre DESC;
\end{lstlisting}

Esta consulta es exactamente la del Ejercicio 3.a del práctico.

% -----------------------------------------------------
\subsection{Reuniones (\texttt{JOIN})}
% -----------------------------------------------------

Los \texttt{JOIN} se escriben dentro del \texttt{FROM} y se corresponden con
las operaciones de reunión del álgebra relacional vista en Bases de Datos I.
Su sintaxis general es:

\begin{lstlisting}
tabla1 [NATURAL] { {LEFT|RIGHT|FULL} OUTER | INNER } JOIN tabla2
       [ON condicionJoin] [USING (columna1 [, columna2, ...])]
\end{lstlisting}

\begin{itemize}
    \item \texttt{INNER JOIN}: devuelve sólo las filas que matchean en ambas
    tablas. Es lo que hace, implícitamente, el viejo estilo
    \texttt{FROM a, b WHERE a.k = b.k}.
    \item \texttt{LEFT OUTER JOIN}: conserva todas las filas de la tabla
    izquierda, completando con \texttt{NULL} cuando no hay coincidencia
    a la derecha. Útil para detectar ``filas huérfanas''.
    \item \texttt{RIGHT OUTER JOIN}: simétrico, conserva todas las de la
    derecha.
    \item \texttt{FULL OUTER JOIN}: conserva todas las filas de ambas
    (no soportado en MySQL nativamente, pero sí en PostgreSQL y Oracle).
    \item \texttt{NATURAL JOIN}: une por todas las columnas con el mismo
    nombre. Cómodo, pero peligroso, porque depende del nombrado.
    \item \texttt{USING (col)}: como el natural pero limitado a las columnas
    explícitas.
\end{itemize}

\paragraph{Ejemplo equivalente.}

\begin{lstlisting}
-- Estilo "moderno" con INNER JOIN.
SELECT p.Nart, pv.Nprov, p.nombre, pv.precio_venta
  FROM proveedores p
  INNER JOIN provee pv ON pv.Nprov = p.Nprov;

-- Estilo "tradicional" con coma y WHERE (equivalente).
SELECT p.Nart, pv.Nprov, p.nombre, pv.precio_venta
  FROM proveedores p, provee pv
 WHERE pv.Nprov = p.Nprov;
\end{lstlisting}

% -----------------------------------------------------
\subsection{Operaciones sobre conjuntos}
% -----------------------------------------------------

SQL permite combinar resultados de varios \texttt{SELECT} con operadores de
conjuntos, siempre y cuando ambas consultas devuelvan la misma cantidad de
columnas y de tipos compatibles:

\begin{itemize}
    \item \texttt{UNION}: unión, eliminando duplicados.
    \item \texttt{UNION ALL}: unión sin eliminar duplicados (más eficiente).
    \item \texttt{INTERSECT}: intersección.
    \item \texttt{EXCEPT} (o \texttt{MINUS} en Oracle): diferencia.
\end{itemize}

% -----------------------------------------------------
\subsection{\texttt{INSERT}}
% -----------------------------------------------------

\texttt{INSERT} agrega filas a una tabla. Tiene dos formas habituales:

\begin{lstlisting}
-- (a) Insertar un registro especificando valores.
INSERT INTO destino (campo1 [, campo2, ...])
VALUES (valor1 [, valor2, ...]);

-- (b) Insertar resultados de una consulta.
INSERT INTO destino (campo1 [, campo2, ...]) seleccion;
\end{lstlisting}

Si no se nombran las columnas, se asume que se está cargando \emph{toda} la
tabla en el orden de definición. Las cadenas se delimitan con comillas
simples y las fechas suelen escribirse en formato \texttt{'YYYY-MM-DD'}.

\paragraph{Ejemplo.}

\begin{lstlisting}
INSERT INTO articulos (Nart, descripcion, precio, stock_max, stock_min)
VALUES (5, 'Vino x1L', 2500, 90, 20);
\end{lstlisting}

% -----------------------------------------------------
\subsection{\texttt{UPDATE}}
% -----------------------------------------------------

\texttt{UPDATE} cambia valores de filas existentes. Su sintaxis es:

\begin{lstlisting}
UPDATE tabla
   SET campo1 = valor1 [, campo2 = valor2, ...]
 [WHERE condicion];
\end{lstlisting}

\textbf{Cuidado:} si se omite el \texttt{WHERE}, la modificación se aplica a
todas las filas de la tabla.

\paragraph{Ejemplo.}

\begin{lstlisting}
-- Aumenta un 20% solo a los articulos baratos.
UPDATE articulos
   SET precio = precio * 1.2
 WHERE precio < 2000;
\end{lstlisting}

% -----------------------------------------------------
\subsection{\texttt{DELETE}}
% -----------------------------------------------------

\texttt{DELETE} elimina filas de una tabla:

\begin{lstlisting}
DELETE FROM tabla [WHERE condicion];
\end{lstlisting}

Sin \texttt{WHERE} elimina \emph{todas} las filas (la tabla queda vacía pero
sigue existiendo, a diferencia de \texttt{DROP TABLE}). Si la fila a borrar
está referenciada por una clave foránea, el resultado depende de la acción
declarada (\texttt{CASCADE}, \texttt{RESTRICT}, \texttt{SET NULL},
\texttt{SET DEFAULT}, \texttt{NO ACTION}); ver Sección~\ref{sec:fk}.

% =====================================================
\section{DDL: definición del esquema}
% =====================================================

% -----------------------------------------------------
\subsection{Tipos de datos}
% -----------------------------------------------------

El estándar SQL clasifica los tipos en 13 categorías primarias. Cada motor
puede agregar variantes y sinónimos. Los tipos del estándar SQL-89 son
\texttt{INTEGER}, \texttt{SMALLINT}, \texttt{CHARACTER}, \texttt{DECIMAL},
\texttt{NUMERIC}, \texttt{REAL}, \texttt{FLOAT} y \texttt{DOUBLE PRECISION};
SQL-92 agrega \texttt{CHARACTER VARYING} (\texttt{VARCHAR}), \texttt{DATE},
\texttt{TIME}, \texttt{BIT}, \texttt{TIMESTAMP}, \texttt{INTERVAL} y
\texttt{BIT VARYING}; SQL:1999 incorpora características objeto-relacionales.

\paragraph{Tipos numéricos (MySQL).}

\begin{center}
\begin{tabular}{@{}l p{8.5cm}@{}}
\toprule
\textbf{Tipo} & \textbf{Descripción} \\
\midrule
\texttt{TINYINT}    & Entero pequeño $-128..127$ (sin signo $0..255$). \\
\texttt{SMALLINT}   & $-32768..32767$. \\
\texttt{MEDIUMINT}  & $-8.388.608..8.388.607$. \\
\texttt{INT}        & $-2.147.483.648..2.147.483.647$. \\
\texttt{BIGINT}     & Hasta $\sim 9.2 \times 10^{18}$. \\
\texttt{FLOAT}      & Coma flotante de precisión simple. \\
\texttt{DOUBLE}     & Coma flotante de precisión doble. \\
\texttt{DECIMAL(p,e)} & Decimal exacto con $p$ dígitos y $e$ decimales. \\
\bottomrule
\end{tabular}
\end{center}

\paragraph{Tipos de fecha y cadena.}

\begin{center}
\begin{tabular}{@{}l p{8.5cm}@{}}
\toprule
\textbf{Tipo} & \textbf{Descripción} \\
\midrule
\texttt{DATE}      & Fecha (\texttt{YYYY-MM-DD}). \\
\texttt{DATETIME}  & Fecha y hora. \\
\texttt{TIMESTAMP} & Fecha y hora con rango acotado, vinculado a la zona
                     horaria del servidor. \\
\texttt{CHAR(n)}   & Cadena de longitud fija. \\
\texttt{VARCHAR(n)}& Cadena de longitud variable hasta $n$. \\
\texttt{TEXT/BLOB} & Textos o binarios extensos. \\
\bottomrule
\end{tabular}
\end{center}

\begin{nota}
\textbf{Para el práctico.} Para precios y montos conviene
\texttt{DECIMAL(p,e)} (evita errores de redondeo de \texttt{FLOAT}); para
identificadores numéricos, \texttt{INT}; para fechas, \texttt{DATE}; para
horarios de entrada/salida (Estacionamiento), \texttt{TIME}.
\end{nota}

% -----------------------------------------------------
\subsection{\texttt{CREATE TABLE}}
% -----------------------------------------------------

\texttt{CREATE TABLE} es el comando central de DDL: define una nueva relación
con sus columnas y sus restricciones. Su forma general es:

\begin{lstlisting}
CREATE TABLE tabla (
   campo1 tipo [(tamano)] [NOT NULL] [restriccion_inline],
   campo2 tipo [(tamano)] [NOT NULL] [restriccion_inline],
   ...
   [, ligaduraIntegridad1]
   [, ligaduraIntegridad2]
);
\end{lstlisting}

Cada campo declara su nombre, tipo y, opcionalmente, si admite \texttt{NULL}
y restricciones a nivel de columna. Al final del bloque pueden agregarse
restricciones de tabla (claves primarias compuestas, foráneas, \texttt{CHECK}s
multi-columna, etc.).

\paragraph{Ejemplo del teórico.}

\begin{lstlisting}
CREATE TABLE Proveedores (
    Nprov INTEGER NOT NULL CONSTRAINT pkproveedores PRIMARY KEY,
    Nombre VARCHAR(50) NOT NULL CONSTRAINT uniconombre UNIQUE,
    Direccion VARCHAR(50)
);
\end{lstlisting}

% -----------------------------------------------------
\subsection{Ligaduras de integridad}
% -----------------------------------------------------
\label{sec:fk}

Las \emph{ligaduras de integridad} (también llamadas restricciones o
\emph{constraints}) garantizan que los datos cumplan ciertas reglas. Las que
maneja SQL para este práctico son:

\begin{itemize}
    \item \textbf{Claves primarias} (\texttt{PRIMARY KEY}): identifican de
    manera única a cada fila. Implican \texttt{NOT NULL} y \texttt{UNIQUE}.
    Una tabla tiene a lo sumo una. Pueden ser simples o \emph{compuestas}.
    \item \textbf{Claves secundarias o candidatas} (\texttt{UNIQUE}):
    aseguran unicidad sobre una o varias columnas, pero permiten
    \texttt{NULL}.
    \item \textbf{No nulidad} (\texttt{NOT NULL}): el campo no admite
    \texttt{NULL}.
    \item \textbf{Restricciones de dominio} (\texttt{CHECK (condicion)}):
    expresan reglas declarativas sobre los valores que un campo puede tomar
    (rangos, predicados, etc.).
    \item \textbf{Integridad referencial} (\texttt{FOREIGN KEY \dots\
    REFERENCES \dots}): obliga a que el valor de una columna exista en otra
    tabla, o sea \texttt{NULL}.
    \item \textbf{Aserciones} (\texttt{ASSERTION}) y \textbf{disparadores}
    (\texttt{TRIGGER}): mecanismos más generales (no se usan en este
    práctico, pero se mencionan en el teórico).
\end{itemize}

La sintaxis general de la cláusula \texttt{CONSTRAINT} es:

\begin{lstlisting}
-- Indices de un solo campo:
CONSTRAINT nombre {PRIMARY KEY | UNIQUE | NOT NULL
                 | REFERENCES tablaexterna [(campoexterno)]}

-- Indices multi-campo:
CONSTRAINT nombre {
    PRIMARY KEY (col1 [, col2, ...])
  | UNIQUE       (col1 [, col2, ...])
  | NOT NULL     (col1 [, col2, ...])
  | CHECK (condicion)
  | FOREIGN KEY (ref1 [, ref2, ...])
        REFERENCES tablaexterna [(extcol1 [, extcol2, ...])]
        [ON DELETE accion] [ON UPDATE accion]
}
\end{lstlisting}

\paragraph{Acciones referenciales.}
Cuando se borra o actualiza una fila a la que apunta una clave foránea, la
acción a tomar puede ser:

\begin{center}
\begin{tabular}{@{}l p{10.5cm}@{}}
\toprule
\textbf{Acción} & \textbf{Comportamiento} \\
\midrule
\texttt{CASCADE}     & Propaga el borrado/actualización a las filas
                       referenciantes. \\
\texttt{RESTRICT}    & Aborta la operación si existen referencias
                       (lo mismo que \texttt{NO ACTION} en muchos motores). \\
\texttt{SET NULL}    & Pone \texttt{NULL} en la columna referenciante. \\
\texttt{SET DEFAULT} & Pone el valor por defecto. \\
\texttt{NO ACTION}   & Variante de \texttt{RESTRICT}, normalmente diferida
                       al final de la transacción. \\
\bottomrule
\end{tabular}
\end{center}

\paragraph{Ejemplos prácticos.}

\begin{lstlisting}
-- Articulo con CHECKs sobre stock y precio.
CREATE TABLE articulos (
    Nart INTEGER NOT NULL PRIMARY KEY,
    Descr VARCHAR(100),
    Precio FLOAT,
    Cant INTEGER,
    Stock_min INTEGER,
    Stock_max INTEGER,
    CONSTRAINT control_min_max CHECK (stock_min < stock_max),
    CONSTRAINT precio_positivo CHECK (precio > 0)
);

-- Provee con clave primaria compuesta y dos claves foraneas.
CREATE TABLE provee (
    Nprov INTEGER NOT NULL,
    Nart INTEGER NOT NULL,
    Precio_venta FLOAT,
    CONSTRAINT pkprovee PRIMARY KEY (Nart, Nprov),
    CONSTRAINT fkproveedores FOREIGN KEY (Nprov) REFERENCES proveedores,
    CONSTRAINT fkarticulos   FOREIGN KEY (Nart)
        REFERENCES articulos ON DELETE CASCADE
);
\end{lstlisting}

\begin{nota}
\textbf{Cómo elegir la acción referencial.} En el inciso 1.a del práctico,
borrar un \emph{cliente} debe arrastrar sus facturas $\Rightarrow$
\texttt{ON DELETE CASCADE}. Pero borrar un \emph{producto} que ya fue
facturado debe ser rechazado $\Rightarrow$ \texttt{ON DELETE RESTRICT}.
La consigna fija el comportamiento esperado y eso determina la acción.
\end{nota}

% -----------------------------------------------------
\subsection{Dominios (\texttt{CREATE DOMAIN})}
% -----------------------------------------------------

Un \emph{dominio} es un tipo definido por el usuario sobre un tipo existente,
opcionalmente con valor por defecto, no nulidad y \texttt{CHECK} asociado:

\begin{lstlisting}
CREATE DOMAIN nombreDominio [AS] tipoDatos
    [DEFAULT valorPorDefecto]
    [NOT NULL]
    [CHECK (condicion)];
\end{lstlisting}

\paragraph{Ejemplo del teórico.}

\begin{lstlisting}
CREATE DOMAIN enteroPositivo AS INTEGER
    DEFAULT 2
    CHECK (VALUE > 1)
    NOT NULL;
\end{lstlisting}

\begin{nota}
\textbf{Soporte por motor.} \texttt{CREATE DOMAIN} es nativo en
\textbf{PostgreSQL}, pero \textbf{MySQL} y \textbf{Oracle} no lo
implementan tal cual el estándar. La cátedra lo emula con la combinación
\emph{tipo de dato + \texttt{CHECK}} (en MySQL también \texttt{ENUM} para
conjuntos cerrados como \texttt{\{gris, negro, azul\}} en el inciso 1.b).
\end{nota}

% -----------------------------------------------------
\subsection{\texttt{ALTER TABLE}}
% -----------------------------------------------------

Permite modificar la estructura de una tabla ya creada:

\begin{lstlisting}
ALTER TABLE tabla {
    ADD {COLUMN campo tipo [(tamano)] [NOT NULL]
       | clausulaConstraint}
  | DROP {COLUMN campo | CONSTRAINT clausulaEliminar}
};
\end{lstlisting}

\paragraph{Ejemplos.}

\begin{lstlisting}
ALTER TABLE Proveedores ADD Telefono VARCHAR(20) NOT NULL;
ALTER TABLE Proveedores DROP Direccion;
ALTER TABLE Proveedores ADD CONSTRAINT telefonounico UNIQUE (Telefono);
ALTER TABLE Proveedores DROP CONSTRAINT telefonounico;
\end{lstlisting}

% -----------------------------------------------------
\subsection{Índices}
% -----------------------------------------------------

Un \emph{índice} es una estructura auxiliar (típicamente B-tree o hash)
basada en una o más columnas, cuya función es acelerar el acceso a los
datos. Los motores lo mantienen automáticamente conforme se insertan,
modifican o borran filas.

\begin{lstlisting}
CREATE [UNIQUE] INDEX indice ON tabla
    (campo1 [ASC|DESC] [, campo2 [ASC|DESC], ...])
    [WITH {PRIMARY | DISALLOW NULL | IGNORE NULL}];
\end{lstlisting}

\textbf{Ideas principales.}
\begin{itemize}
    \item \texttt{UNIQUE INDEX}: además de acelerar, impide duplicados.
    \item \texttt{PRIMARY KEY} y \texttt{UNIQUE} crean índices implícitamente.
    \item Cada \texttt{INSERT}/\texttt{UPDATE}/\texttt{DELETE} cuesta más con
    cada índice; los índices aceleran consultas pero penalizan modificaciones.
\end{itemize}

% -----------------------------------------------------
\subsection{Vistas (\texttt{CREATE VIEW})}
% -----------------------------------------------------

Una \emph{vista} es una relación virtual, definida como una consulta
nominada. No forma parte del modelo lógico, pero se ve y se usa como una
tabla más:

\begin{lstlisting}
CREATE VIEW <nombre_vista> AS <expresion_de_consulta>;

-- Ejemplo.
CREATE VIEW nombre_proveedores AS
    SELECT nombre FROM proveedores;
\end{lstlisting}

Sus usos clásicos son: ocultar columnas sensibles (seguridad), simplificar
consultas frecuentes y desacoplar el esquema lógico de las aplicaciones que
lo consultan. Su \emph{actualización} (\texttt{UPDATE} sobre una vista) es
problemática y depende de cuán mapeable sea la vista a una sola tabla base.

% -----------------------------------------------------
\subsection{Borrado de estructuras: \texttt{DROP}}
% -----------------------------------------------------

\texttt{DROP} elimina objetos definitivamente. Su sintaxis básica es:

\begin{lstlisting}
DROP {TABLE tabla | INDEX indice ON tabla};
\end{lstlisting}

A diferencia de \texttt{DELETE}, que vacía una tabla pero la conserva,
\texttt{DROP TABLE} la elimina completamente, incluyendo su definición.

% =====================================================
\section{Conexión con los ejercicios del práctico}
% =====================================================

Los tres ejercicios del práctico recorren los temas anteriores en el orden
natural: primero el esquema (DDL), después los datos (\texttt{INSERT}) y
finalmente las consultas (\texttt{SELECT} con todas sus cláusulas).

\subsection*{Ejercicio 1 --- Definición de esquemas}

Tres bases de datos independientes (\texttt{practico1a}, \texttt{practico1b},
\texttt{practico1c}) que ejercitan distintos rincones de DDL:

\begin{itemize}
    \item \textbf{Inciso a} (MySQL y PostgreSQL). Cliente-Producto-Factura-
    \mbox{ItemFactura}. Foco: \texttt{PRIMARY KEY} simple y compuesta,
    \texttt{CHECK} (\texttt{precio > 0}, \texttt{stock\_minimo <=
    stock\_maximo}) y \texttt{FOREIGN KEY} con acciones distintas
    (\texttt{ON DELETE CASCADE} para facturas de un cliente borrado,
    \texttt{ON DELETE RESTRICT} para productos ya facturados).
    \item \textbf{Inciso b} (MySQL y Oracle). Vehículo-Persona-Dueño-
    Estacionamiento-Parquímetro. Foco: \texttt{CREATE DOMAIN}
    (\emph{nombreYApellido} como \texttt{VARCHAR(45)}), restricción de
    conjunto cerrado (\emph{color} en \{gris, negro, azul\} con
    \texttt{ENUM} en MySQL o \texttt{CHECK} en Oracle), \texttt{CHECK} de
    rango (\texttt{altura BETWEEN 0 AND 5000}), autoincrementales
    (\texttt{AUTO\_INCREMENT} en MySQL, \texttt{SEQUENCE + TRIGGER} o
    \texttt{IDENTITY} en Oracle) y \texttt{ON DELETE RESTRICT} para vehículos
    estacionados.
    \item \textbf{Inciso c} (PostgreSQL y Oracle). Cliente-Automóvil-
    Categoría-Taller-Accidente. Foco: dominio de marcas restringido
    (\{FIAT, RENAULT, FORD\}), rango sobre el modelo (1990–2015), control
    sobre la tarifa (positiva y menor a 10000) y \texttt{ON DELETE CASCADE}
    para los accidentes ante el borrado de un cliente.
\end{itemize}

\subsection*{Ejercicio 2 --- Carga de datos coherentes}

Para cada una de las tres bases, se insertan datos consistentes que respeten
las claves primarias, las foráneas y los \texttt{CHECK}. ``Coherente'' aquí
implica además que los datos alcancen para que las consultas del
Ejercicio 3 produzcan resultados no triviales: si pediremos clientes sin
facturas, debe existir alguno; si pediremos clientes con más de tres
accidentes, también.

\subsection*{Ejercicio 3 --- Consultas}

Cada inciso ejercita un patrón distinto de \texttt{SELECT}:

\begin{itemize}
    \item \textbf{a)} ``Clientes sin facturas'' es un caso típico de
    \emph{antijoin}: se resuelve con \texttt{NOT EXISTS} o \texttt{NOT IN}
    (o \texttt{LEFT JOIN \dots\ WHERE \dots\ IS NULL}). Se cierra con un
    \texttt{ORDER BY apellido DESC, nombre DESC}.
    \item \textbf{b)} ``Vehículos que usaron el parquímetro 9'' requiere un
    \texttt{JOIN} entre \emph{vehiculo} y \emph{estacionamiento} más un
    filtro sobre \texttt{nro\_parquimetro = 9}. Se devuelven sólo
    \emph{modelo} y \emph{color}, casi siempre con \texttt{DISTINCT}
    para evitar repeticiones por cada estacionamiento.
    \item \textbf{c)} ``Clientes con más de 3 accidentes'' es el caso clásico
    de \texttt{GROUP BY} + \texttt{HAVING}: se agrupa por DNI, se cuenta y
    se filtra con \texttt{HAVING COUNT(*) > 3}.
    \item \textbf{d)} ``Máximo y mínimo monto de factura por cliente'' usa
    \texttt{MAX}, \texttt{MIN} y \texttt{GROUP BY nro\_cliente}, junto con un
    \texttt{JOIN} hacia \emph{cliente} para mostrar el nombre.
    \item \textbf{e)} Tres consultas en álgebra relacional, donde aparezcan
    todos los operadores: \emph{selección} ($\sigma$), \emph{proyección}
    ($\pi$), \emph{unión} ($\cup$), \emph{intersección} ($\cap$) y
    \emph{producto cartesiano} ($\times$). Es el puente con la materia
    Bases de Datos I.
\end{itemize}

\paragraph{Ejemplo resuelto (3.c).}

\begin{lstlisting}
SELECT c.dni, c.nombre, c.apellido
  FROM cliente c
  JOIN accidente a ON a.dni = c.dni
 GROUP BY c.dni, c.nombre, c.apellido
HAVING COUNT(*) > 3;
\end{lstlisting}

\paragraph{Ejemplo resuelto (3.d).}

\begin{lstlisting}
SELECT c.nro_cliente,
       c.nombre,
       MAX(f.monto) AS monto_max,
       MIN(f.monto) AS monto_min
  FROM cliente c
  JOIN factura f ON f.nro_cliente = c.nro_cliente
 GROUP BY c.nro_cliente, c.nombre;
\end{lstlisting}

% =====================================================
\section{Diferencias entre motores: lo que no hay que olvidar}
% =====================================================

\begin{center}
\begin{tabular}{@{}p{3.4cm} p{3.6cm} p{3.6cm} p{3.6cm}@{}}
\toprule
\textbf{Tema} & \textbf{MySQL 8+} & \textbf{PostgreSQL 14+} &
\textbf{Oracle 10/11 XE} \\
\midrule
``Crear la base'' & \texttt{CREATE DATABASE} &
\texttt{CREATE DATABASE} (cada una corre independiente) &
Se crea un \emph{usuario/esquema}: \texttt{CREATE USER \dots} y se
trabaja sobre él. \\
\addlinespace
Auto incremento & \texttt{AUTO\_INCREMENT} en columna &
\texttt{SERIAL} / \texttt{GENERATED \dots\ AS IDENTITY} &
\texttt{SEQUENCE} + \texttt{TRIGGER} (10/11) o \texttt{IDENTITY} (12+). \\
\addlinespace
Dominios & No soporta \texttt{CREATE DOMAIN}; emular con
\texttt{CHECK} o \texttt{ENUM} & Soporta \texttt{CREATE DOMAIN} estándar &
No soporta \texttt{CREATE DOMAIN}; emular con \texttt{CHECK}. \\
\addlinespace
Conjuntos cerrados & \texttt{ENUM('a','b','c')} &
\texttt{CHECK (col IN ('a','b','c'))} o \texttt{CREATE TYPE \dots\ AS ENUM} &
\texttt{CHECK (col IN ('a','b','c'))}. \\
\addlinespace
Booleano & \texttt{BOOLEAN} (sinónimo de \texttt{TINYINT(1)}) &
\texttt{BOOLEAN} real & No tiene; emular con \texttt{NUMBER(1)} y
\texttt{CHECK}. \\
\addlinespace
Texto largo & \texttt{TEXT}, \texttt{LONGTEXT} & \texttt{TEXT} &
\texttt{CLOB}. \\
\bottomrule
\end{tabular}
\end{center}

% =====================================================
\section{Ideas clave para repasar antes del parcial}
% =====================================================

\begin{enumerate}
    \item \textbf{SQL se divide en DML, DDL y DCL.} Saber rápidamente a qué
    grupo pertenece cada sentencia ahorra confusiones.
    \item \textbf{La sentencia \texttt{SELECT} tiene seis cláusulas} y un
    \emph{orden lógico} de evaluación distinto al orden de escritura.
    \texttt{FROM} $\to$ \texttt{WHERE} $\to$ \texttt{GROUP BY} $\to$
    \texttt{HAVING} $\to$ \texttt{SELECT} $\to$ \texttt{ORDER BY}.
    \item \textbf{\texttt{WHERE} filtra filas; \texttt{HAVING} filtra grupos.}
    Si no hay agregadas, no hace falta \texttt{HAVING}.
    \item \textbf{Las funciones de agregación ignoran \texttt{NULL}}, salvo
    \texttt{COUNT(*)}.
    \item \textbf{\texttt{NULL} se compara con \texttt{IS NULL} / \texttt{IS
    NOT NULL}}, nunca con \texttt{=}.
    \item \textbf{\texttt{LIKE} usa \texttt{\%} y \texttt{\_}} como
    comodines para patrones de cadenas.
    \item \textbf{\texttt{JOIN} es la versión moderna del producto cartesiano
    + \texttt{WHERE}}, y los \texttt{OUTER JOIN} son la forma natural de
    encontrar ``filas sin pareja''.
    \item \textbf{Una clave primaria es \texttt{NOT NULL} + \texttt{UNIQUE}}.
    Las claves candidatas se modelan con \texttt{UNIQUE}; las foráneas con
    \texttt{REFERENCES}.
    \item \textbf{Las acciones referenciales} (\texttt{CASCADE},
    \texttt{RESTRICT}, \texttt{SET NULL}, \texttt{SET DEFAULT},
    \texttt{NO ACTION}) son las que bajan los requisitos del enunciado al
    SQL.
    \item \textbf{\texttt{CREATE DOMAIN} es estándar y de PostgreSQL}; en
    MySQL y Oracle se emula con \texttt{CHECK} o \texttt{ENUM}. Ese punto
    es el que diferencia los tres incisos del Ejercicio 1.
    \item \textbf{\texttt{INSERT}, \texttt{UPDATE}, \texttt{DELETE} sin
    \texttt{WHERE}} actúan sobre toda la tabla. Es uno de los errores más
    comunes y costosos.
    \item \textbf{\texttt{DROP TABLE} elimina la tabla completa};
    \texttt{DELETE FROM tabla} sólo borra las filas.
    \item \textbf{Patrones típicos de consulta del práctico}:
    \begin{itemize}
        \item ``A sin B'' $\Rightarrow$ \texttt{NOT EXISTS} o \texttt{LEFT
        JOIN \dots\ WHERE \dots\ IS NULL}.
        \item ``Cuántos B por A'' $\Rightarrow$ \texttt{GROUP BY a} con
        \texttt{COUNT(*)}.
        \item ``A con más de $k$ B'' $\Rightarrow$ \texttt{GROUP BY a
        HAVING COUNT(*) > $k$}.
        \item ``Máx/mín por A'' $\Rightarrow$ \texttt{GROUP BY a} con
        \texttt{MAX/MIN}.
    \end{itemize}
    \item \textbf{Cada motor tiene sus modismos.} La diferencia más
    importante a la hora de codificar es la creación del esquema/base, los
    autoincrementales y el manejo de dominios y enums.
\end{enumerate}

% =====================================================
\section*{Resumen final}
% =====================================================
\addcontentsline{toc}{section}{Resumen final}

El Práctico 1 es un repaso integral que reúne casi todos los pilares de SQL:
DDL para definir tablas con claves, restricciones de dominio y de integridad
referencial; DML para poblarlas y consultarlas; \texttt{JOIN}s, agrupamientos
y funciones agregadas para consultas no triviales; y una práctica concreta de
las diferencias entre motores. Si entran sólidas las cinco ideas de
``\texttt{SELECT} con sus seis cláusulas'', ``orden lógico de evaluación'',
``\texttt{WHERE} vs \texttt{HAVING}'', ``acciones referenciales'' y
``dominios estándar vs emulados'', queda muy buena base para encarar
el resto de la materia y, en particular, el Práctico 2 (DCL, permisos y
roles), que se construye exactamente sobre estos esquemas.

\end{document}
